\documentclass{subfiles}
\begin{document}
    Most of the groups and ring will be using in these notes,
        rise naturally from the notion of congruence modulo \(n\),
        where \(n\) is an integer. These groups and rings will be denoted by \(\zset_{n}\),
        or by \(\zset_{p}\) if \(n = p\) is prime.

    But what is congruence? Put in simple words, given two integers \(a, b\),
        we say that \(a \text{ is congruent } b\) modulo \(n\), 
        if we can write \(a\) as the sum of \(b\) and some multiple of \(n\).
        That is, 
        \[
            a \equiv b \mod n \iff \exists k \in \zset : a = b + kn.
        \]
        Using the above definition, we get
        \[
            \zset_{n} = \set{m \in \nset}[m < n].
        \]
        For instance, \(\zset_{11} = \set{0, 1, \ldots, 9, 10}\).
        Sometimes when dealing with rings we will write \(\zset_{n}^{*}\) to mean the group/ring
        \(\zset_{n} \setminus \set{0}\).

    In most of the algorithms we will discuss,
        we often make use of the \(\gcd\) between two given integers.
        Here we first recall what GCD is, and provide two algorithm to compute it.

    Let \(a, b\) be two integers.
        We say that an integer \(d\) is the GCD of \(a, b\) if the following hold:
        \begin{equation}\label{eq:gcd-definition}
            d \vert a \text{ and } d \vert b 
        \end{equation}
        and for any other integer \(d'\) such that \Cref{eq:gcd-definition} holds,
        we have
        \[
            d' \vert d.
        \]
        In other words, the \(d\) is the GCD of \(a, b\) if any integer that divides \(a, b\)
        also divides \(d\).

    But how do we compute the GCD of two integers? A result in abstract algebra tells us
        that we can compute \(d\) as the difference between some multiple of \(a\)
        and some multiple of \(b\). This result is known as \emph{Bézout identity}.
        More precisely Bézout identity states that there exists two integers \(\delta \text{ and } \epsilon\),
        such that \(d = \gcd(a, b) = a\delta + b\epsilon\).
        \(\delta \text{ and } \epsilon\) are called Bézout coefficients.

    In practise we compute the GCD by means of the Euclidean algorithm.
        Two variants exists, a simple one which returns just the GCD,
        and one which returns both the GCD and the Bézout coefficients.
        Here we present both, although for our needs the latter is preferred.

        \subfile{../TikZ/Algorithm 1 - Simple Euclidean algorithm}
        \subfile{../TikZ/Algorithm 2 - Extended Euclidean algorithm}
\end{document}
